% ============================================================
% Colours, boxes, and links
% ============================================================
\usepackage{xcolor}
\usepackage{mdframed}
\usepackage{tcolorbox}
\tcbuselibrary{skins,breakable}

% ============================================================
% TikZ
% ============================================================
\usepackage{tikz}
\usepackage{tikz-3dplot}
\usetikzlibrary{
  arrows.meta,
  backgrounds,
  calc,
  decorations.pathreplacing,
  decorations.pathmorphing,
  decorations.markings,
  fit,
  patterns,
  positioning,
  shapes.geometric
}

% ============================================================
% Miscellaneous
% ============================================================
\usepackage{etoolbox}
\usepackage{expl3}
\usepackage{float}
\usepackage{microtype}
\usepackage{multicol}
\usepackage{placeins}
\usepackage{xparse}
\usepackage{multirow}
\usepackage{rotating}   % sidewaysfigure

% ============================================================
% Bibliography
% ============================================================
\usepackage[authoryear]{natbib}



\usepackage{imakeidx}
\makeindex[name=tech,title={Technique Index}]
\makeindex[name=dep,title={Dependency Index}]


% ============================================================
% exercise-tags-and-stubs-v2.tex
% Include file for the Learning Abstract Mathematics project.
%
% PURPOSE
%   - Defines searchable technique/dependency tag helpers.
%   - Defines sourced exercise-record stubs:
%       * exerciseconstruction      (Types 4, 5)
%       * exerciseproofrecord       (Types 1, 2, 3, 6)
%   - Defines generated exercise-record stubs:
%       * generatedexerciseconstruction
%       * generatedexerciseproof
%   - Defines a quick generated exercise macro:
%       * \exstubgenfull
%
% DESIGN GOALS
%   - No \usepackage calls here.
%   - No theorem-environment redefinitions.
%   - Safe to \input into main.tex after your existing preamble.
%   - Uses only packages already present in your main.tex:
%       xparse, etoolbox, expl3, tcolorbox, hyperref
%
% SEARCHABILITY
%   Write literal tags in source, e.g.
%       tech:triangle-inequality
%       tech:epsilon-splitting
%       def:metric
%       thm:triangle-inequality
%   so that grep/search works on the .tex source directly.
%
% OPTIONAL INDEXING
%   If you later load makeidx/imakeidx and define \index,
%   \TechniqueEntry{...} and \DependencyEntry{...} will add index entries.
%   Otherwise they simply typeset the tag.
% ============================================================

\makeatletter

% ------------------------------------------------------------
% Optional index hooks
% ------------------------------------------------------------
\providecommand{\ExerciseTechniqueIndexName}{tech}
\providecommand{\ExerciseDependencyIndexName}{dep}

\newcommand{\exercise@maybeindextech}[1]{%
  \@ifundefined{index}{}{%
    \index[tech]{#1}%
  }%
}
\newcommand{\exercise@maybeindexdep}[1]{%
  \@ifundefined{index}{}{%
    \index[dep]{#1}%
  }%
}

% ------------------------------------------------------------
% Tag entry macros
%   IMPORTANT: pass the FULL literal tag string.
%   Example:
%     \TechniqueEntry{tech:triangle-inequality}
%     \DependencyEntry{def:metric}
% ------------------------------------------------------------
\providecommand{\TechniqueEntry}[1]{%
  \texttt{#1}%
  \exercise@maybeindextech{#1}%
}

\providecommand{\DependencyEntry}[1]{%
  \texttt{#1}%
  \exercise@maybeindexdep{#1}%
}

% ------------------------------------------------------------
% Convenience wrappers for common tag namespaces
% ------------------------------------------------------------
\providecommand{\tech}[1]{\TechniqueEntry{tech:#1}}
\providecommand{\dep}[1]{\DependencyEntry{#1}}

% ------------------------------------------------------------
% Metadata storage (sourced exercises)
% ------------------------------------------------------------
\providecommand{\CurrentExerciseAuthor}{}
\providecommand{\CurrentExerciseBook}{}
\providecommand{\CurrentExerciseID}{}
\providecommand{\CurrentExerciseType}{}
\providecommand{\CurrentExerciseTechniques}{}
\providecommand{\CurrentExerciseDependencies}{}

% ------------------------------------------------------------
% Metadata storage (generated exercises)
% ------------------------------------------------------------
\providecommand{\CurrentGeneratedExerciseID}{}
\providecommand{\CurrentGeneratedExerciseTopic}{}
\providecommand{\CurrentGeneratedExerciseType}{}
\providecommand{\CurrentGeneratedExerciseTechniques}{}
\providecommand{\CurrentGeneratedExerciseDependencies}{}

% ------------------------------------------------------------
% Per-author counters
% ------------------------------------------------------------
\newcounter{exerciseAbbott}
\newcounter{exerciseTao}
\newcounter{exerciseLebl}
\newcounter{exerciseBruckner}
\newcounter{exerciseBartleSherbert}
\newcounter{exerciseGeneric}

\newcommand{\exercise@stepauthorcounter}[1]{%
  \ifdefstrequal{#1}{Abbott}{\stepcounter{exerciseAbbott}}{%
  \ifdefstrequal{#1}{Tao}{\stepcounter{exerciseTao}}{%
  \ifdefstrequal{#1}{Lebl}{\stepcounter{exerciseLebl}}{%
  \ifdefstrequal{#1}{Bruckner}{\stepcounter{exerciseBruckner}}{%
  \ifdefstrequal{#1}{Bartle \& Sherbert}{\stepcounter{exerciseBartleSherbert}}{%
  \stepcounter{exerciseGeneric}}}}}}%
}

\newcommand{\exercise@authorcountvalue}[1]{%
  \ifdefstrequal{#1}{Abbott}{\arabic{exerciseAbbott}}{%
  \ifdefstrequal{#1}{Tao}{\arabic{exerciseTao}}{%
  \ifdefstrequal{#1}{Lebl}{\arabic{exerciseLebl}}{%
  \ifdefstrequal{#1}{Bruckner}{\arabic{exerciseBruckner}}{%
  \ifdefstrequal{#1}{Bartle \& Sherbert}{\arabic{exerciseBartleSherbert}}{%
  \arabic{exerciseGeneric}}}}}}%
}

% ------------------------------------------------------------
% Metadata setter (sourced exercises)
% ------------------------------------------------------------
\NewDocumentCommand{\ExerciseMeta}{m m m m m m}{%
  % #1 author
  % #2 book
  % #3 exercise id
  % #4 type
  % #5 techniques   (literal source-searchable tags)
  % #6 dependencies (literal source-searchable tags)
  \renewcommand{\CurrentExerciseAuthor}{#1}%
  \renewcommand{\CurrentExerciseBook}{#2}%
  \renewcommand{\CurrentExerciseID}{#3}%
  \renewcommand{\CurrentExerciseType}{#4}%
  \renewcommand{\CurrentExerciseTechniques}{#5}%
  \renewcommand{\CurrentExerciseDependencies}{#6}%
  \exercise@stepauthorcounter{#1}%
}

% ------------------------------------------------------------
% Metadata setter (generated exercises)
% ------------------------------------------------------------
\NewDocumentCommand{\GeneratedExerciseMeta}{m m m m m}{%
  % #1 exercise id
  % #2 topic
  % #3 type
  % #4 techniques
  % #5 dependencies
  \renewcommand{\CurrentGeneratedExerciseID}{#1}%
  \renewcommand{\CurrentGeneratedExerciseTopic}{#2}%
  \renewcommand{\CurrentGeneratedExerciseType}{#3}%
  \renewcommand{\CurrentGeneratedExerciseTechniques}{#4}%
  \renewcommand{\CurrentGeneratedExerciseDependencies}{#5}%
}

% ------------------------------------------------------------
% Shared header (sourced exercises)
% ------------------------------------------------------------
\newcommand{\ExerciseRecordHeader}{%
\begin{tcolorbox}[colback=gray!6,colframe=gray!40,arc=2pt,
  left=6pt,right=6pt,top=4pt,bottom=4pt,
  title={\small\textbf{Exercise Record}}]
\small
\textbf{Author index number. }
\CurrentExerciseAuthor~\exercise@authorcountvalue{\CurrentExerciseAuthor}

\medskip
\textbf{Source. }
\CurrentExerciseAuthor, \textit{\CurrentExerciseBook}, \CurrentExerciseID

\medskip
\textbf{Type. }
\CurrentExerciseType

\medskip
\textbf{Technique tags. }
\texttt{\CurrentExerciseTechniques}

\medskip
\textbf{Dependency tags. }
\texttt{\CurrentExerciseDependencies}
\end{tcolorbox}
}

% ------------------------------------------------------------
% Shared header (generated exercises)
% ------------------------------------------------------------
\newcommand{\GeneratedExerciseRecordHeader}{%
\begin{tcolorbox}[colback=gray!6,colframe=gray!40,arc=2pt,
  left=6pt,right=6pt,top=4pt,bottom=4pt,
  title={\small\textbf{Exercise Record}}]
\small
\textbf{Origin. } Generated exercise

\medskip
\textbf{Exercise ID. } \CurrentGeneratedExerciseID

\medskip
\textbf{Topic. } \CurrentGeneratedExerciseTopic

\medskip
\textbf{Type. } \CurrentGeneratedExerciseType

\medskip
\textbf{Technique tags. } \texttt{\CurrentGeneratedExerciseTechniques}

\medskip
\textbf{Dependency tags. } \texttt{\CurrentGeneratedExerciseDependencies}
\end{tcolorbox}
}

% ------------------------------------------------------------
% Stub A — Construction / Find / Compute (sourced)
% ------------------------------------------------------------
\NewDocumentEnvironment{exerciseconstruction}{}{%
  \subsection*{Exercise}
  \ExerciseRecordHeader

  % Suggested searchable metadata block:
  % tech:...
  % def:...
  % thm:...

  \begin{remark*}[Return]
  $\leftarrow$ Back-link to the relevant notes section.
  \end{remark*}

  \begin{remark*}[Task]
  State the exercise verbatim here.
  \end{remark*}

  \begin{remark*}[Construction]
  State the object or computed value explicitly before any justification.
  \end{remark*}

  \begin{proof}
}{%
  \end{proof}

  \begin{remark*}[Dependencies]
  Expand the dependency tags into prose if helpful.
  \end{remark*}

  \begin{remark*}[Proof shape]
  Record the structural pattern used:
  direct / contradiction / verification / construction / computation.
  \end{remark*}

  \begin{remark*}[Reflection]
  Optional: explain why the construction works and what feature is essential.
  \end{remark*}
}

% ------------------------------------------------------------
% Stub B — Proof / Verify / Prove-or-Counterexample / Characterization (sourced)
% ------------------------------------------------------------
\NewDocumentEnvironment{exerciseproofrecord}{}{%
  \subsection*{Exercise}
  \ExerciseRecordHeader

  % Suggested searchable metadata block:
  % tech:...
  % def:...
  % thm:...

  \begin{remark*}[Return]
  $\leftarrow$ Back-link to the relevant notes section.
  \end{remark*}

  \begin{remark*}[Task]
  State the exercise verbatim here.
  \end{remark*}

  \begin{remark*}[Strategy]
  Briefly describe the proof plan.
  \end{remark*}

  \begin{proof}
}{%
  \end{proof}

  \begin{remark*}[Dependencies]
  Expand the dependency tags into prose if helpful.
  \end{remark*}

  \begin{remark*}[Proof shape]
  Record the structural pattern:
  direct / contradiction / case split / subsequence extraction / squeeze / iff.
  \end{remark*}
}

% ------------------------------------------------------------
% Stub A — Construction / Find / Compute (generated)
% ------------------------------------------------------------
\NewDocumentEnvironment{generatedexerciseconstruction}{}{%
  \subsection*{Exercise}
  \GeneratedExerciseRecordHeader

  % Suggested searchable metadata block:
  % tech:...
  % def:...
  % thm:...

  \begin{remark*}[Task]
  State the generated exercise here.
  \end{remark*}

  \begin{remark*}[Construction]
  State the object or computed value explicitly before justification.
  \end{remark*}

  \begin{proof}
}{%
  \end{proof}

  \begin{remark*}[Dependencies]
  Expand the dependency tags into prose if helpful.
  \end{remark*}

  \begin{remark*}[Proof shape]
  Record the structural pattern used:
  direct / contradiction / verification / construction / computation.
  \end{remark*}

  \begin{remark*}[Reflection]
  Optional: explain why the construction works.
  \end{remark*}
}

% ------------------------------------------------------------
% Stub B — Proof / Verify / Prove-or-Counterexample / Characterization (generated)
% ------------------------------------------------------------
\NewDocumentEnvironment{generatedexerciseproof}{}{%
  \subsection*{Exercise}
  \GeneratedExerciseRecordHeader

  % Suggested searchable metadata block:
  % tech:...
  % def:...
  % thm:...

  \begin{remark*}[Task]
  State the generated exercise here.
  \end{remark*}

  \begin{remark*}[Strategy]
  Briefly describe the proof plan.
  \end{remark*}

  \begin{proof}
}{%
  \end{proof}

  \begin{remark*}[Dependencies]
  Expand the dependency tags into prose if helpful.
  \end{remark*}

  \begin{remark*}[Proof shape]
  Record the structural pattern:
  direct / contradiction / case split / squeeze / iff.
  \end{remark*}
}

% ------------------------------------------------------------
% Quick unsolved generated exercise stub
% ------------------------------------------------------------
\newcommand{\exstubgenfull}[5]{%
  \GeneratedExerciseMeta{#1}{#2}{GENERATED}{#3}{#4}%

  \subsection*{Exercise}
  \GeneratedExerciseRecordHeader

  \begin{remark*}[Task]
  #5
  \end{remark*}

  \begin{remark*}[Work Area]
  Write your proof or computation here.
  \end{remark*}

  \bigskip
}

% ------------------------------------------------------------
% Controlled vocabularies
% ------------------------------------------------------------

% TECHNIQUE TAGS (recommended)
% tech:unpack-definition
% tech:triangle-inequality
% tech:epsilon-splitting
% tech:tail-argument
% tech:choose-N
% tech:case-splitting
% tech:contradiction
% tech:subsequence-extraction
% tech:diagonal-argument
% tech:squeeze
% tech:monotone-bound
% tech:supremum-characterization
% tech:metric-verification
% tech:ball-containment
% tech:diameter-bound
% tech:compactness-argument

% DEPENDENCY TAGS (recommended)
% def:metric
% def:open-ball
% def:bounded
% def:diameter
% def:convergence
% def:cauchy-sequence
% def:supremum
% def:infimum
% thm:triangle-inequality
% thm:archimedean-property
% thm:monotone-convergence
% thm:bolzano-weierstrass
% thm:squeeze-theorem

% ------------------------------------------------------------
% Example usage (sourced)
% ------------------------------------------------------------
%
% % =========================================================
% % Exercise Metadata
% % tech:triangle-inequality
% % tech:epsilon-splitting
% % def:convergence
% % thm:triangle-inequality
% % =========================================================
%
% \ExerciseMeta
%   {Abbott}
%   {Understanding Analysis}
%   {Exercise 2.3.8}
%   {PROOF}
%   {tech:triangle-inequality, tech:epsilon-splitting}
%   {def:convergence, thm:triangle-inequality}
%
% \begin{exerciseproofrecord}
%
% \begin{remark*}[Technique]
% \TechniqueEntry{tech:triangle-inequality},
% \TechniqueEntry{tech:epsilon-splitting}
% \end{remark*}
%
% \begin{remark*}[Dependencies]
% \DependencyEntry{def:convergence},
% \DependencyEntry{thm:triangle-inequality}
% \end{remark*}
%
% Let $\varepsilon > 0$. Since $x_n \to x$, there exists $N$ such that
% $d(x_n,x)<\varepsilon/2$ for all $n\ge N$. If $m,n\ge N$, then
% \[
% d(x_m,x_n)\le d(x_m,x)+d(x,x_n)<\varepsilon.
% \]
%
% \end{exerciseproofrecord}
%
% ------------------------------------------------------------
% Example usage (generated)
% ------------------------------------------------------------
%
% % =========================================================
% % Exercise Metadata
% % tech:supremum-characterization
% % tech:unpack-definition
% % def:supremum
% % def:infimum
% % =========================================================
%
% \exstubgenfull
%   {EX-BOUNDS-CHATGPT-001}
%   {Bounds / Suprema / Infima}
%   {tech:supremum-characterization, tech:unpack-definition}
%   {def:supremum, def:infimum}
%   {Let $A=\{1/n : n\in\mathbb N\}$. Determine $\sup A$ and $\inf A$.}
%
% ------------------------------------------------------------

\makeatother





% --- Blueprint macros (merged from blueprint-macros.tex) ---
% ============================================================
%  Blueprint shared macros — \input{blueprint-macros} once
%  after \begin{document} before any blueprint \input calls.
%
%  Required in main.tex preamble:
%    \usepackage{tcolorbox}
%    \tcbuselibrary{skins}
%    \usepackage{amsthm,amsmath,amssymb,xcolor}
%    all \definecolor{cSet}...{cVS} definitions
%    \newtheorem{definition}{Definition}[chapter]  (or similar)
%    \newtheorem{remark}{Remark}[chapter]
% ============================================================

% ── Section boxes ─────────────────────────────────────────────
% NEW: strong frame in structure colour
\newtcolorbox{bpnew}[2]{
  enhanced,
  colback   = #1!14!white,
  colframe  = #1!80!black,
  coltitle  = white,
  fonttitle = \bfseries\scshape\small,
  title     = #2,
  boxrule   = 1.0pt, arc=2pt,
  left=7pt, right=7pt, top=4pt, bottom=5pt,
  toptitle=3pt, bottomtitle=3pt,
  before skip=0pt, after skip=7pt,
}

% INHERITED: faded frame in ancestor colour
\newtcolorbox{bpinh}[2]{
  enhanced,
  colback   = #1!7!white,
  colframe  = #1!38!white,
  coltitle  = #1!70!black,
  fonttitle = \bfseries\scshape\small,
  title     = #2,
  boxrule   = 0.6pt, arc=2pt,
  left=7pt, right=7pt, top=4pt, bottom=5pt,
  toptitle=3pt, bottomtitle=3pt,
  before skip=0pt, after skip=7pt,
}

% ── Entry line ────────────────────────────────────────────────
% \bpitem{symbol}{math}{gloss}   (gloss may be empty)
\newcommand{\bpitem}[3]{%
  \noindent
  {\small\bfseries #1}\hspace{6pt}%
  $\displaystyle #2$%
  \ifx\relax#3\relax\else
    \hfill{\small\rmfamily\itshape\color{lbltext!75}\,[#3]}%
  \fi
  \par\vspace{3pt}%
}

% ── Page header ───────────────────────────────────────────────
% \bpheader{colour}{Name}{signature}
\newcommand{\bpheader}[3]{%
  \noindent\vbox{\offinterlineskip
    \hbox{\colorbox{#1!85!black}{%
      \parbox[c][2.4em]{\dimexpr\linewidth-2\fboxsep\relax}{%
        \color{nodetext}%
        \hspace{8pt}{\large\bfseries #2}%
        \enspace{\normalsize\ttfamily\color{sigtext} #3}%
        \hfill{\footnotesize\scshape\bfseries Structure Blueprint\hspace{8pt}}%
      }%
    }}%
  }\par\vspace{8pt}%
}

% ── Footer ────────────────────────────────────────────────────
% \bpfooter{left}{right}
\newcommand{\bpfooter}[2]{%
  \vfill
  \noindent\hrule\vspace{4pt}%
  \noindent
  \begin{minipage}[t]{0.49\linewidth}\scriptsize #1\end{minipage}%
  \hfill
  \begin{minipage}[t]{0.49\linewidth}\scriptsize #2\end{minipage}%
}

% =========================================================
% Dependency Figure Styles
% Add to common/boxes.tex
% =========================================================

% Requires TikZ libraries:
% \usetikzlibrary{arrows.meta,positioning,calc}

\tikzset{
  depaxiom/.style={
    draw,
    rounded corners=3pt,
    fill=axiombox,
    draw=axiomborder,
    minimum width=4.2cm,
    minimum height=0.82cm,
    align=center,
    font=\small
  },
  depdefinition/.style={
    draw,
    rounded corners=3pt,
    fill=defbox,
    draw=defborder,
    minimum width=4.2cm,
    minimum height=0.82cm,
    align=center,
    font=\small
  },
  deptheorem/.style={
    draw,
    rounded corners=3pt,
    fill=thmbox,
    draw=thmborder,
    minimum width=4.2cm,
    minimum height=0.82cm,
    align=center,
    font=\small
  },
  depneutral/.style={
    draw,
    rounded corners=3pt,
    fill=gray!8,
    draw=gray!65,
    minimum width=4.2cm,
    minimum height=0.82cm,
    align=center,
    font=\small
  },
  depimplies/.style={->, thick, blue!70!black, >=Latex},
  depuses/.style={->, thick, yellow!60!black, >=Latex},
  depdefines/.style={->, thick, green!50!black, >=Latex},
depedgelab/.style={
  font=\tiny,
  fill=white,
  inner sep=0.5pt,
  text=black
},
  deplegendbox/.style={
    draw,
    rounded corners=3pt,
    fill=white,
    draw=gray!55,
    inner sep=5pt
  }
}

\newcommand{\DependencyLegend}{%
\begin{tikzpicture}[baseline=(current bounding box.north)]
\node[deplegendbox, anchor=north east] (L) {
\begin{tabular}{@{}ll@{}}
\multicolumn{2}{@{}l@{}}{\textbf{\small Legend}}\\[0.2em]
\tikz{\node[depaxiom, minimum width=0.7cm, minimum height=0.28cm] {};} & \scriptsize Axiom\\
\tikz{\node[depdefinition, minimum width=0.9cm, minimum height=0.35cm] {};} & \scriptsize Definition\\
\tikz{\node[deptheorem, minimum width=0.9cm, minimum height=0.35cm] {};} & \scriptsize Theorem / Proposition / Corollary\\
\tikz{\draw[depimplies] (0,0) -- (0.8,0);} & \scriptsize implies\\
\tikz{\draw[depuses] (0,0) -- (0.8,0);} & \scriptsize uses\\
\tikz{\draw[depdefines] (0,0) -- (0.8,0);} & \scriptsize defines
\end{tabular}
};
\end{tikzpicture}%
}
% =========================================================
% Dependency Figure Legend
% Standard bottom placement legend for dependency maps
% =========================================================

\newcommand{\DependencyFigureLegend}{%
\begin{tikzpicture}[baseline=(current bounding box.center), >=Latex]
  % Row 1: node legend
  \node[depaxiom, minimum width=0.9cm, minimum height=0.32cm] (ax) at (0,0) {};
  \node[right=0.15cm of ax] (axlab) {\scriptsize Axiom};

  \node[depdefinition, minimum width=0.9cm, minimum height=0.32cm, right=0.9cm of axlab] (df) {};
  \node[right=0.15cm of df] (dflab) {\scriptsize Definition};

  \node[deptheorem, minimum width=0.9cm, minimum height=0.32cm, right=0.9cm of dflab] (th) {};
  \node[right=0.15cm of th] (thlab) {\scriptsize Theorem / Proposition / Corollary};

  % Row 2: edge legend
  \draw[depimplies] ($(ax.south)+(0,-0.9)$) -- ++(0.9,0)
    node[right, text=black] {\scriptsize implies};

  \draw[depuses] ($(df.south)+(0,-0.9)$) -- ++(0.9,0)
    node[right, text=black] {\scriptsize uses};

  \draw[depdefines] ($(th.south)+(0,-0.9)$) -- ++(0.9,0)
    node[right, text=black] {\scriptsize defines};
\end{tikzpicture}%
}
% ---------------------------------------------------------
% Minimal template for a dependency figure file
% Save as figure-<n>.tex and \input it from the notes.
% ---------------------------------------------------------
%
% \begin{center}
% \textbf{<Subchapter or Chapter Title> Dependency Map}
% \end{center}
%
% \begin{tikzpicture}[>=Latex]
%   % Nodes
%   \node[depaxiom]      (a1) at (0,0) {Completeness Axiom};
%   \node[deptheorem]    (t1) at (-3,-1.7) {Nested Interval\\Property};
%   \node[deptheorem]    (t2) at ( 3,-1.7) {Existence of $\sqrt{a}$};
%   \node[deptheorem]    (t3) at (-3,-3.4) {Archimedean Property};
%
%   % Edges
%   \draw[depimplies] (a1) -- node[depedgelab, pos=0.55, left] {implies} (t1);
%   \draw[depimplies] (a1) -- node[depedgelab, pos=0.55, right] {implies} (t2);
%   \draw[depuses]    (t1) -- node[depedgelab, pos=0.50, right] {uses} (t3);
%
%   % Legend (place manually where it fits)
%   \node[anchor=north east] at (6.0,0.2) {\DependencyLegend};
% \end{tikzpicture}

% =========================================================
% Topic-level pedagogical containers
% =========================================================

\newtcolorbox{topicbox}[1]{
  enhanced,
  breakable,
  colback=white,
  colframe=white,
  arc=2pt,
  boxrule=0pt,
  title={\textbf{#1}},
  fonttitle=\bfseries,
  coltitle=black,
  colbacktitle=white,
  boxed title style={boxrule=0pt, colframe=white, colback=white},
  attach boxed title to top left={xshift=0pt, yshift*=-\tcboxedtitleheight/2},
  frame hidden,
  left=0pt, right=0pt, top=8pt, bottom=4pt,
  before skip=8pt, after skip=8pt
}

\NewDocumentEnvironment{exposition}{}%
  {\par\smallskip\noindent\ignorespaces}%
  {\par\smallskip}

\newcommand{\NoLocalDependencies}{}


% ============================================================
% Formal mathematical content boxes
% ============================================================

\newtcolorbox{definitionbox}[1]{
  colback=defbox,
  colframe=defborder,
  arc=2pt,
  left=6pt,
  right=6pt,
  top=4pt,
  bottom=4pt,
  title={\small\textbf{#1}},
  fonttitle=\small\bfseries
}

\newtcolorbox{axiombox}[1]{
  colback=axiombox,
  colframe=axiomborder,
  arc=2pt,
  left=6pt,
  right=6pt,
  top=4pt,
  bottom=4pt,
  title={\small\textbf{#1}},
  fonttitle=\small\bfseries
}

\newtcolorbox{theorembox}[1]{
  colback=thmbox,
  colframe=thmborder,
  arc=2pt,
  left=6pt,
  right=6pt,
  top=4pt,
  bottom=4pt,
  title={\small\textbf{#1}},
  fonttitle=\small\bfseries
}

\newtcolorbox{lemmabox}[1]{
  colback=lembox,
  colframe=lemborder,
  arc=2pt,
  left=6pt,
  right=6pt,
  top=4pt,
  bottom=4pt,
  title={\small\textbf{#1}},
  fonttitle=\small\bfseries
}

\newtcolorbox{propositionbox}[1]{
  colback=propbox,
  colframe=propborder,
  arc=2pt,
  left=6pt,
  right=6pt,
  top=4pt,
  bottom=4pt,
  title={\small\textbf{#1}},
  fonttitle=\small\bfseries
}

\newtcolorbox{corollarybox}[1]{
  colback=corbox,
  colframe=corborder,
  arc=2pt,
  left=6pt,
  right=6pt,
  top=4pt,
  bottom=4pt,
  title={\small\textbf{#1}},
  fonttitle=\small\bfseries
}


% ============================================================
% House colors: navigation / breadcrumbs
% ============================================================
\definecolor{breadcrumb}{HTML}{F2F2F2}
\definecolor{breadcrumbborder}{HTML}{B8B8B8}
\definecolor{breadcrumbtitle}{HTML}{333333}

\providecolor{amber}{HTML}{F59E0B}
